% 第 47 章　自旋：没有经典对应物的角动量（完整章节）
\chapter{自旋：没有经典对应物的角动量}

前面几章一直在给量子世界立新规矩：状态用振幅记账（第~43~章），两次测量之间按薛定谔方程演化（第~44~章），测量是一次会改写状态的真实相互作用（第~45~章），而这些账目可以整齐地写成向量和矩阵（第~46~章）。本章要请出量子力学里最出名、也最被误解的一个角色：自旋。它是一种角动量，却不对应任何东西的转动；它只有两种测量结果，却撑起了磁铁的磁性、医院的磁共振成像和量子计算机的比特。学完本章，你应该能亲手解释一块冰箱磁铁为什么吸得住铁钉，也能说清楚为什么任何“电子是一颗自转小球”的图画都注定要失败。

\step{反常识开场}

花几块钱网购一块钕磁铁（俗称强磁，冰箱贴、耳机、手机扬声器里都有），把它靠近一串回形针：它能隔着空气把第一枚回形针拽过来，第一枚又拽起第二枚，像拎起一条铁链子。现在换一块同样大小的铜块或者铝块，无论怎么凑近，回形针纹丝不动。

这件事你从小看到大，习惯了，但请认真问一句：凭什么？铜和铁都由原子构成，原子核外都拖着一群带负电的电子。第~33~章讲过，运动的电荷产生磁场——电子在原子里时刻处于运动状态，按理说每个原子都该是个小磁体，铜块里的电子并不比铁块里的懒。那为什么一个磁、一个不磁？磁性到底藏在哪里？

再往下问一层，事情更怪。把条形磁铁从中间锯开，你不会得到“一个北极块”和一个“南极块”，而是两块完整的小磁铁，各自带南北两极。再锯，还是如此。磁性似乎长在每个原子的骨子里，可原子又是电中性的。一条合理的猜测是：每个电子本身就是一枚微型磁针。这个猜测是对的——但紧接着的下一个猜测，“那它一定是一颗在自转的带电小球”，将把我们带进一整章的麻烦。磁性这条线，最终会牵出一种你在经典力学里从未见过的角动量：它不转，却真实存在；它让银原子束在磁铁里劈成两半；它此刻正在你身体里的每个氢原子核上，等待一台磁共振仪来点名。

\step{先猜一猜}

老规矩，先把你的直觉写在纸上，越具体越好。

第一题。假设电子真的是一颗带电的小球，绕着过球心的轴自转，就像一颗微型地球。带电物质转起来形成环形电流，环形电流产生磁矩——这幅图像解释磁性似乎天经地义。你猜：要转出实验测得的电子磁性，这颗小球的表面线速度大概是多少？是每秒几米（像陀螺），每秒几千米（像炮弹），还是每秒几万千米（接近光速）？请先写下一个数，稍后我们会真的算它。

第二题。把一束银原子加热蒸发出来，让它们通过一对磁极——磁极做成一尖一平的形状，使磁场沿竖直方向不均匀地变化。原子如果带磁矩，就会被不均匀的磁场推拉：小磁针“北极朝上”的被推向一边，“南极朝上”的被推向另一边。你猜屏幕上会留下什么痕迹？（a）一条连续展宽的竖直带子，因为原子的磁针朝向应该什么角度都有；（b）上下两个分离的斑点；（c）一个斑点，因为原子太小根本推不动。写下你的选择和理由。

第三题。假设实验给出了（b），两个斑点。现在把“被推到上方”的那群原子单独挑出来，再让它们通过一台\textbf{同样朝向}的磁极装置，你猜它们还会分开吗？然后把装置原地转九十度，让磁场沿水平方向变化，这群“上方原子”会怎样？最后，把转过的装置再转回来，重新测一次竖直方向，结果还和第一次筛选之后一样吗？这一串问题听起来繁琐，但它们的答案正是本章最锋利的那把刀。

\step{动手实验}

自旋本身无法在家里看到——这需要真空、高温炉和精密磁铁。但本章要和一种经典角动量反复对照，所以请先把“普通的转动”看个真切，知道我们要告别的是什么。

\begin{experiment}[转椅与指南针：先认清经典的角动量和磁针]
\textbf{材料}：一把能顺畅旋转的办公椅（或圆凳），两个灌满水的矿泉水瓶当哑铃；外加一枚指南针（手机的指南针应用也可以，但要远离磁铁校准）。

\textbf{第一步}：坐在转椅上，双臂平举，两手各握一个水瓶。让同伴轻轻推你转起来，然后双脚离地。把双臂猛地收拢到胸前——转速明显变快；再张开，转速变慢。这就是角动量守恒：转动惯量变小，角速度必须变大来补偿。花样滑冰运动员收臂加速，用的全是这一条。

\textbf{第二步}：在旋转中注意一件事——角动量是有方向的（沿转轴）。试着在旋转时把整个身体连同椅子倾斜（小心扶稳），你会感到椅子轴“不愿意”歪，有股劲在抵抗。方向是角动量的核心属性，不只是转多快，还有“绕哪根轴、朝哪边转”。

\textbf{第三步}：把指南针放在桌上，记下它指向。拿钕磁铁从不同方位靠近：磁针转动，最终停在顺着磁铁磁场方向的姿态；把磁铁掉头，磁针也跟着掉头。磁针在磁场里像一个总想把自家电磁矩对准外场方向的罗盘。注意：磁针的指向是\textbf{连续可调}的——你把外场转到任意角度，它就跟到任意角度，中间没有任何被跳过的方向。

\textbf{记录与思考}：写下两条结论。其一，经典角动量的大小和方向都可以连续变化，你想让它多快、朝哪都行。其二，磁针在外磁场中可以稳定地指向任何方向。请记住这两条——它们是“旧模型”的全部本钱，本章稍后要证明：微观世界里的那种角动量，这两条一条都不满足。
\end{experiment}

如果找不到转椅，单做第三步也有收获：用两块磁铁互相靠近，感受“只有两个取向彼此最舒服”（同向排列相吸、首尾相接），但任何中间角度也都允许存在——只是需要你的手一直拧着。连续的取向，连续的力矩，这就是经典磁针。

\step{旧模型遇到麻烦}

现在来算那道你猜过的题：电子如果是一颗自转的带电小球，它要转多快？

实验测得电子携带的角动量大小是 $\hbar/2$，其中 $\hbar\approx1.05\times10^{-34}$~J·s，所以电子的“自转角动量”约为 $5.3\times10^{-35}$~J·s。经典力学里，一个转动物体的角动量大致等于质量乘以半径再乘以表面线速度（差一个与形状有关、数量级为一的系数）。电子的质量是 $9.1\times10^{-31}$~kg。电子的半径取多少？如果电子有内部结构，高能对撞实验早该看到端倪，但直到目前最深的探测，电子仍然表现得像一个点，其半径上限远小于 $10^{-18}$~m。我们\textbf{故意放宽}，取经典电子半径 $r\approx2.8\times10^{-15}$~m——这是按经典电磁图像能给电子安上的最大块的“身材”。于是
\[
v=\frac{L}{m\,r}\approx\frac{5.3\times10^{-35}}{9.1\times10^{-31}\times2.8\times10^{-15}}
\approx2\times10^{10}\ \text{m/s}.
\]
光速是 $3\times10^{8}$~m/s。也就是说，哪怕把电子的身材放到最宽，它表面的“转动速度”也要达到光速的\textbf{六七十倍}；如果采用实验给出的半径上限，这个数还要再膨胀上千倍。狭义相对论（第~38~章）严禁任何物质运动超过光速。结论很硬：电子的角动量不可能来自任何“自转”。这不是技术细节上的小出入，而是六七十倍起步的、原则性的破产。顺便检查这个估算的方向：把半径取得越小，所需速度越大，麻烦越严重——所以“也许电子其实更小”救不了这个模型，只会让它死得更透。

第二个麻烦来自那束银原子。1922 年，斯特恩和盖拉赫真的做了第二题里的实验：银原子从加热炉里出来，经过狭缝整成细束，穿过一尖一平的磁极之间那段不均匀的磁场，打到玻璃屏上。先想清楚经典预言：原子的磁针（如果每个原子是一枚小磁针的话）从火炉里出来时朝向应该完全随机，上下左右什么角度都有。不均匀的磁场对磁针的推力正比于磁矩沿磁场方向的分量，朝向随机的磁针，分量也连续地随机分布，于是屏上的斑点应该被连续地涂抹开来，形成一条竖直的带子——最上端是磁针完全顺场的原子，最下端是完全逆场的，中间密密麻麻排满各种斜着的。这就是经典的、也是几乎所有物理学家当年的预期。

实验结果：玻璃屏上是\textbf{两个}分离的斑点，一上一下，中间干干净净，什么都没有。没有带子，没有过渡，仿佛自然界只允许这枚磁针有两种朝向。你猜对了吗？当年连斯特恩和盖拉赫自己都被这个结果惊到——后来人们查明，银原子的磁性几乎全部来自它最外层那颗孤独的电子，所以屏上的劈裂实际上是电子磁矩、也就是电子自旋的劈裂。

第三个麻烦更阴险，就是第三题那串追问。把上方斑点的原子挑出来（用挡板把下方那束挡住），送进第二台同方向的装置：它们\textbf{全部}再次向上偏，不再分裂。看来第一次筛选确实挑出了一种确定的属性，叫它“上”——到此为止还很经典，像用筛子筛出了大颗的豆子，再过同样的筛子当然还是大颗。接下来把第二台装置转九十度，让磁场沿水平方向不均匀：这束“上”原子\textbf{又劈成了两束}，左一半右一半。也还不算太离奇——磁针嘛，换个方向测，竖直的磁针对水平分量确实可正可负。真正的惊雷在最后一步：把“左”这一束挑出来，送入第三台转回竖直方向的装置。按照一切经典经验，这束原子的祖辈明明被认证过是“上”原子，竖直属性从未被任何人改动过，第三台装置应该全部给出“上”。实验结果：\textbf{一半上，一半下}。那次水平测量不像是“读出”了一个既有的水平属性，倒像是把竖直方向的档案\textbf{销毁重写}了。筛豆子模型死无全尸——筛过的大豆绝不会因为被称了一次重量就变出一半小豆子来。

这幅“测量之间互相擦掉对方结果”的图景，你在第~45~章的三片偏振片实验里已经见过一次：竖直偏振的光子通过四十五度偏振片后，水平分量被重新“造”了出来。银原子束的分裂与重组，是同一套量子逻辑在角动量上的重演。这不是巧合，这是在提示：我们面对的不是某个实验装置的怪癖，而是“状态”这个概念本身的新规则。

\step{发明新概念}

账目烂掉了，就重新设计账簿。我们一步步来，每一步都只做实验逼我们做的事。

第一步，给那种劈不开的角动量起名字。电子携带一份内禀的角动量，大小固定为 $\hbar/2$，与它的位置、速度、是否束缚在原子里全部无关——它就是电子身份的一部分，像质量和电荷一样。我们叫它\textbf{自旋}。请注意这个名字是个历史事故：它暗示转动，而我们刚刚证明了它不能是转动。物理学家也曾想把电子换成更大的复合球来挽救图像，全部失败；今天最深刻的理论（第~49~章的狄拉克方程）显示，自旋是把量子力学和狭义相对论缝在一起时\textbf{自动掉出来}的东西，根本不是谁在里面旋转。所以请把“自旋”当作一个专名来用，就像“鲸鱼”不是鱼，自旋也不是旋。

第二步，承认它的测量结果少得反常。沿任意选定的方向去测电子自旋，结果只有两种：沿该方向“向上”（记作 $+\hbar/2$）或“向下”（记作 $-\hbar/2$）。注意这个“任意”有多蛮横：你选竖直方向测，只有上、下；改选水平方向测，还是只有左、右两个结果；斜着四十三度测，依然只有两个结果。经典的磁针可以有连续的分量——与磁场夹角三十度时，沿场分量就是总磁矩乘以 $\cos30^\circ$——电子的自旋却像一枚只有两面、却能沿任何轴抛掷的硬币。这就是为什么我们称电子为\textbf{二能级系统}：就一个方向的提问而言，它的全部状态空间只有两格。这个“两格”的系统还有一个更现代的名字：\textbf{量子比特}。经典比特只能站在 0 或 1 上；量子比特的状态可以是这两个基态的任意叠加——这正是一部分量子计算机（第~48~章会再谈）的物理地基。

第三步，规定状态之间如何换算。第~45~章的偏振已经教会我们规则：选定一个方向，就有了“上”“下”两个基态；对另一个方向提问时，把旧状态按新方向的两个基态展开，展开系数的模平方给出两个结果的概率，测量之后状态更新为对应的新基态。自旋完全照此办理。于是一切怪异自动归位：第一次竖直测量把原子制备成“竖直上”态；水平装置问的是一个新问题，“竖直上”态相对水平方向展开成“左”“右”各占一半，所以第二台装置劈成两束；测得“左”的原子此刻处于“水平左”态，这个态相对竖直方向又是上、下各半——第三台装置给出一半上一半下。连锁测量里“前一次的档案被后一次销毁”，不是装置的粗鲁，而是二能级系统状态几何的直接推论。

第四步，把自旋和磁性焊在一起。电子的自旋带着一个磁矩，大小约为一个玻尔磁子 $\mu_{\mathrm B}\approx9.3\times10^{-24}$~J/T，方向与自旋绑定（电子带负电，所以磁矩与自旋方向相反，这个细节只影响符号，不影响本章任何结论）。银原子束在磁场里受力，本质就是这份磁矩在不均匀磁场里被推拉：磁矩顺场的原子被拉向磁场强的一侧，逆场的被推向弱的一侧，于是两个斑点对应磁矩——也就是自旋——的两种取向。自此，磁铁的磁性、磁共振成像的信号、磁盘的存储，全部可以追溯到这个没有经典对应物的小小角动量。

值得立刻给一个比喻，再立刻给它拴上缰绳。自旋像什么？像一枚硬币：抛一次只有正反两面，你可以换不同的姿势抛（沿不同的轴测），但每一抛仍然只有正反两个结果。它不像什么？不像的地方恰恰是硬币的全部常识：硬币抛之前，正反已经写在硬币的朝向上，测量只是抄答案；而电子在被问到某个方向之前，并不携带着沿该方向的确定答案——第~45~章的教训在这里原样生效，而且第~48~章的贝尔实验将证明，连“答案预先写好、只是没读”这种隐变量式补救在数学上都走不通。还要声明一条：说电子“像磁针”，指的仅仅是它的磁矩在外场中有两种投影结果；它不像真实磁针的地方在于，真实磁针的磁矩来自电荷的宏观环流、取向连续，电子的磁矩却是内禀属性、测量结果离散。

\step{一句话模型}

\begin{oneline}
自旋是粒子与生俱来的一份角动量：它不来自任何转动，大小永远固定（电子为 $\hbar/2$），沿任何方向测量都只有“顺”“逆”两种结果；换方向测量就是换一组新问题，旧方向的答案随之作废——而这份小小的角动量绑着磁矩，是世间一切磁性与磁共振的总开关。
\end{oneline}

这句话像什么？像户口：自旋是电子一出生就登记在册的属性，不依赖它后来经历了什么。它不像什么？不像体检报告：体检指标会因你的活动而改变，而 $\hbar/2$ 这个值对全宇宙每一颗电子分毫不差——你能“转动”的只是提问的方向。再补一条“它不是”：自旋不是电子还没学会的经典角动量的迷你版；恰恰相反，你手上陀螺的角动量，才是亿万个自旋与轨道运动叠加之后的宏观近似。

\step{数学翻译器}

先给二能级系统一套记号。选定一个方向（叫它 $z$ 方向），把两个基态写成 $|\,{\uparrow}_z\rangle$ 和 $|\,{\downarrow}_z\rangle$。一个沿与 $z$ 轴夹角 $\theta$ 方向“向上”的自旋态，按 $z$ 方向展开为
\[
|\,{\uparrow}_\theta\rangle=\cos\frac{\theta}{2}\;|\,{\uparrow}_z\rangle+\sin\frac{\theta}{2}\;|\,{\downarrow}_z\rangle .
\]
按老规矩问四个问题。它描述什么：同一个自旋状态，用另一组测量基来记账。为什么这样写：系数必须满足模平方和为一（概率总归一），且要保证 $\theta=0$ 时纯上、$\theta=180^\circ$ 时纯下——半角形式是自旋特有的记账方式，来源超出了主线，挑战层盒子会给一个入口。何时成立：理想、无相互作用的单次投影测量。何时失效：真实装置有退磁、有退相干、有探测器效率，它只是这些效应都被压低后的极限。

玻恩规则立刻给出本章最核心的一个数：一束“$\theta$ 方向上”的原子，通过沿 $z$ 方向的斯特恩—盖拉赫装置时，出现在上方通道的概率是
\[
P_{\uparrow}=\cos^{2}\frac{\theta}{2}.
\]
立刻用特殊情况盘问它。$\theta=0$：$\cos^2 0=1$，全部向上——对，同一台装置同方向复测不该分裂。$\theta=180^\circ$：$\cos^2 90^\circ=0$，全部向下——对，方向完全反转，结果完全反转。$\theta=90^\circ$：$\cos^2 45^\circ=\tfrac12$，水平态在竖直装置里劈成两半——这正是连锁实验第二台装置看到的一比一。注意一个细节：自旋的概率里是\textbf{半角}，而第~45~章偏振光子里是 $\cos^2\theta$ 的整角。这个“二分之一”是自旋独有的胎记，它在挑战层会酿成一件著名的事：自旋态转 $360^\circ$ 不回到原样。

再看期望值。对大量相同的 $\theta$ 态原子测 $z$ 方向自旋，平均结果（以 $\hbar/2$ 为单位）是
\[
\langle S_z\rangle=\frac{\hbar}{2}\cos^{2}\frac{\theta}{2}-\frac{\hbar}{2}\sin^{2}\frac{\theta}{2}
=\frac{\hbar}{2}\cos\theta .
\]
检查两端：$\theta=0$ 时平均就是 $+\hbar/2$，$\theta=180^\circ$ 时是 $-\hbar/2$，$\theta=90^\circ$ 时平均为零——单次结果从来只有 $\pm\hbar/2$，“$\cos\theta$”那个连续值只活在\textbf{大量重复的平均}里。请把这条刻进脑子：微观上离散，宏观平均上连续；你在转椅实验里摸到的那个连续可调的角动量世界，就是这样从离散的底层平均出来的。

\begin{figure}[htbp]
\centering
\begin{tikzpicture}[>=Stealth,scale=1.0]
  % 炉子
  \draw[thick] (0,-0.55) rectangle (1.0,0.55);
  \node at (0.5,0) {银炉};
  \draw[thick,->] (1.0,0) -- (2.0,0);
  % 磁极
  \draw[thick,fill=gray!25] (2.4,0.9) -- (3.6,0.9) -- (3.2,0.35) -- (2.8,0.35) -- cycle;
  \node[above] at (3.0,0.95) {N};
  \draw[thick,fill=gray!25] (2.4,-0.9) rectangle (3.6,-0.5);
  \node[below] at (3.0,-0.95) {S};
  % 分裂轨迹
  \draw[thick] (2.0,0) -- (2.6,0);
  \draw[thick] (3.6,0.15) .. controls (4.6,0.35) .. (5.8,0.55);
  \draw[thick] (3.6,-0.15) .. controls (4.6,-0.35) .. (5.8,-0.55);
  % 屏
  \draw[thick] (5.9,-1.0) -- (5.9,1.0);
  \fill (5.9,0.55) circle (0.06);
  \fill (5.9,-0.55) circle (0.06);
  \node[right] at (5.95,0) {屏：两个斑点};
\end{tikzpicture}
\caption{斯特恩—盖拉赫装置示意。不均匀磁场（尖极 N、平极 S）把银原子束按自旋取向劈成两束，屏上只有两个分离斑点，而不是经典预言的连续带。}
\end{figure}

现在把磁矩算成真实数字，看看那两个斑点需要多大的力气。设磁场沿竖直方向的不均匀程度（梯度）为每米变化 $10$~T，这是当年装置容易达到的量级。磁矩为 $\mu_{\mathrm B}$ 的原子在不均匀磁场中受的力约为 $F=\mu_{\mathrm B}\times(\text{梯度})\approx9.3\times10^{-24}\times10\approx9\times10^{-23}$~N。银原子质量约 $1.8\times10^{-25}$~kg，于是加速度 $a=F/m\approx500$~m/s$^2$——是重力加速度的五十倍。原子以每秒几百米的热速度飞过约 $0.1$~m 长的磁极区，用时约 $2\times10^{-4}$~s，横向偏移 $\tfrac12 a t^{2}\approx10^{-5}$~m，即十微米量级，两束分开几十微米。当年玻璃屏上那两个肉眼可辨的小斑点，背后就是这么一笔账：十的负二十三次方牛顿的力，推在十的负二十五次方千克的原子上，积少成多，在宏观世界里盖了章。

\begin{figure}[htbp]
\centering
\begin{tikzpicture}[>=Stealth,scale=1.0]
  \draw[thick,->] (-0.4,0) -- (10.6,0) node[right]{原子束};
  % 三台装置
  \foreach \x/\lab in {1.6/{$z$ 装置}, 5.2/{转过 $90^\circ$}, 8.8/{转回 $z$}}{
    \draw[thick,fill=blue!8] (\x-0.45,-0.8) rectangle ++(0.9,1.6);
    \node[below] at (\x,-0.85) {\lab};
  }
  \node[above] at (0.6,0.1) {未筛选};
  % 分支
  \draw[thick] (2.05,0.35) -- (3.1,0.6) node[right]{上：留};
  \draw[thick,dashed] (2.05,-0.35) -- (3.1,-0.6) node[right]{下：挡住};
  \draw[thick] (3.35,0.6) -- (4.75,0.6);
  \draw[thick] (5.65,0.75) -- (6.7,1.05) node[right]{左 $\tfrac12$：留};
  \draw[thick,dashed] (5.65,0.45) -- (6.7,0.15) node[right]{右 $\tfrac12$：挡住};
  \draw[thick] (6.95,1.05) -- (8.35,1.05);
  \draw[thick] (9.25,1.2) -- (10.2,1.5) node[right]{上 $\tfrac12$};
  \draw[thick] (9.25,0.9) -- (10.2,0.6) node[right]{下 $\tfrac12$？！};
\end{tikzpicture}
\caption{三台斯特恩—盖拉赫装置串联。第一次筛出的“上”原子，被水平方向测量“洗”过之后，再次测量竖直方向竟重新出现一半“下”——前次筛选的档案被后次测量销毁。}
\end{figure}

\begin{mathbridge}
\textbf{泡利矩阵：二能级系统的全部代数。}把 $|\,{\uparrow}_z\rangle$ 和 $|\,{\downarrow}_z\rangle$ 写成二维列向量 $\binom{1}{0}$ 与 $\binom{0}{1}$，沿三个方向测自旋（以 $\hbar/2$ 为单位）对应三个矩阵
\[
\sigma_x=\begin{pmatrix}0&1\\[2pt]1&0\end{pmatrix},\qquad
\sigma_y=\begin{pmatrix}0&-i\\[2pt]i&0\end{pmatrix},\qquad
\sigma_z=\begin{pmatrix}1&0\\[2pt]0&-1\end{pmatrix}.
\]
几何直觉这样建立：$\sigma_z$ 的本征向量正是两个基向量，本征值 $\pm1$——沿 $z$ 测，结果非上即下，矩阵对角。$\sigma_x$ 的本征向量是 $\tfrac{1}{\sqrt2}\binom{1}{\pm1}$，即“上”“下”的等权叠加——所以 $z$ 态用 $x$ 装置测，五五开；把上式反过来，$x$ 本征态用 $z$ 装置测同样五五开。矩阵之间不对易：$\sigma_x\sigma_z\neq\sigma_z\sigma_x$，差值正比于 $\sigma_y$。第~46~章说过，不对易就是“测量顺序会改变账”，连锁实验里那次“档案销毁”，在代数上就是这三个矩阵互相乘不对易。一般方向的自旋态写作
\[
|\,{\uparrow}_{\theta,\varphi}\rangle=\cos\frac{\theta}{2}\,|\,{\uparrow}_z\rangle+e^{i\varphi}\sin\frac{\theta}{2}\,|\,{\downarrow}_z\rangle ,
\]
其中 $\theta,\varphi$ 就是普通球坐标里的两个角：每一个自旋态对应球面上的一个点，这个球叫布洛赫球。量子比特的全部状态，就是这颗球的整个球面——经典比特只有南北两极，量子比特把整颗球都给了你。
\end{mathbridge}

\begin{challenge}
\textbf{转两圈才回家：旋量与 $720^\circ$。}半角系数带来一个著名后果：把自旋态的参数 $\theta$ 之外的整体转动角推进 $360^\circ$，$\cos(\theta/2)$ 与 $\sin(\theta/2)$ 各自变号，态向量变成原来的 $-1$ 倍；要转 $720^\circ$ 才回到自身。单个态的这个 $-1$ 不改变任何测量概率（概率是模平方），所以它不是印刷错误，而是一类新对象的身份证：这种随转动按半角变换的量叫\textbf{旋量}。$360^\circ$ 与 $720^\circ$ 的差别甚至可以宏观演示：手掌向上托平一本书，保持书不翻转，手臂原地绕身体转一圈，你的胳膊会拧住；用同样的约束再转一圈，拧住的胳膊反而解开了——拓扑上，转两圈的路径可以连续收缩回不动，转一圈不行（这叫狄拉克的“盘子戏法”）。数学背景是：三维旋转群与自旋变换群之间是二对一的覆盖关系。还有一个更深的预告：电子磁矩与自旋的比值（$g$ 因子）在狄拉克方程里自然等于 $2$（第~49~章），量子电动力学把它修正到 $2.002319\ldots$，与实验吻合到小数点后十几位——这是物理学史上最精确的预言之一，而它的起点，就是本章这个“不转动的角动量”。
\end{challenge}

\step{模型的边界}

第一，自旋不是转动，相关的一切直觉都要设防。“电子像一颗旋转的星球”这个画面会立刻要你回答“赤道线速度多少”，而我们已经算过，答案是光速的几十倍起步——这个模型在算术上自杀。安全的说法是：自旋是一个在转动操作下按特定规则（旋量规则）变换的内禀量子数，它\textbf{贡献}角动量、参与角动量守恒，却不对应任何空间中的环流。角动量守恒依然成立：电子自旋翻转时，角动量差额会精确地转移给光子或周围环境，账从没烂过。

第二，“自旋向上”永远要带上方向。说“这个电子自旋向上”而不指明沿哪个轴，就像说“这座山高两千”而不说是海拔还是相对山脚。而且沿 $z$ 方向确定的态，沿 $x$ 方向必然是完全不确定的叠加——不是信息不足，是态本身就是这样。这不是仪器精度问题，仪器做得再完美，单次结果依然只有 $\pm\hbar/2$、比例依然五五开。

第三，半整数与整数的分界线。电子、质子、中子的自旋是 $\hbar/2$（半整数），光子是 $\hbar$（整数）。这条分界不是分类癖：半整数自旋的粒子遵守泡利不相容原理——两个电子不能占据完全相同的量子态，这正是原子有壳层结构、元素周期表有周期、你脚下的地板是固体的深层原因；整数自旋的粒子则相反，可以任意聚集在同一状态，激光的相干光束就是例子。第~53~章前后的统计物理章节会展开这条线，这里只先立一块路标：本章那颗小小的 $\hbar/2$，最终决定了物质为什么不塌陷、不透明、有体积。

第四，磁共振成像测的是\textbf{群体}，不是单个自旋。医院里没有任何仪器在追踪你某一颗质子的自旋；磁共振信号来自亿万个质子磁矩在磁场中的\textbf{平均磁化}的进动与弛豫。单次测量只有两种结果，宏观信号却是连续的波形——这正是“微观离散、平均连续”那条规则的工程版。把磁共振想象成“读出每颗质子的状态”会立刻撞上不确定关系；把它理解为“对宏观磁化矢量的射频激励与监听”，一切就顺理成章。

第五，别给自旋乱安意识或通信功能。自旋翻转不传递任何信息，单个自旋的测量结果完全随机；两个自旋可以纠缠（第~48~章），但纠缠同样不能用来超光速传信。市面上一切“量子共振检测健康”“自旋能量水”之类的产品，与本章物理的关系是零。

\step{回头解释开场现象}

现在回到那串回形针。钕磁铁为什么磁，铜块为什么不磁？

每个电子都是一枚内禀小磁针，这一点铜和铁、钕一视同仁。差别出在\textbf{配对}与\textbf{排队}两件事上。在绝大多数材料里，电子两两配对占据同一个轨道，而泡利不相容原理要求配对的两个电子自旋相反——两枚小磁针首尾相抵，磁性干干净净地抵消掉。铜的所有电子都处在这种配对状态，所以铜块对磁场几乎无动于衷。而钕、铁这类元素的原子里有\textbf{没配上对}的电子；更关键的是，在这些固体中，相邻原子的未配对电子之间存在一种量子力学特有的交换相互作用（它本质上是泡利原理加静电排斥的联合推论，不是磁力也不是谁“指挥”谁），使得邻居的自旋\textbf{平行排列}能量更低。亿万个自旋于是排成同方向的大军，每颗贡献一个 $\mu_{\mathrm B}$，宏观叠加出能吊起回形针串的磁场。铁被钕磁铁吸引，则是因为铁里的自旋大军在外场诱导下重新整队——这就是“磁化”。至于加热到一定温度（居里温度，铁约 $770\,^\circ$C）磁性消失，则是热运动把排队打散、阵列瓦解的故事，那是热学章节的领地。

锯开磁铁为什么总是得到两块完整磁铁？因为磁性不在磁铁的“两端”，而在每一个自旋上；锯子只是把已经排好队的自旋大军切成两支队伍，每支队伍内部仍然整齐。世上没有磁单极子可被锯出来——至少在本书写作时，实验上从未找到。

磁共振成像则是同一套物理最辉煌的应用。你的身体约六成是水，每个水分子携带两个氢原子核——质子，质子也是自旋 $\hbar/2$ 的粒子，自带磁矩。把人放进磁共振仪的强磁场（常用 $1.5$~T，是地球磁场的三万倍），质子自旋只有两个取向：顺场与逆场，两种取向的能量差 $\Delta E$ 正比于磁场强度。换算成频率，质子磁矩在 $1.5$~T 磁场中的进动频率约为 $64$~MHz——恰好落在无线电广播的波段。仪器于是发射这个频率的射频脉冲，像给秋千按节奏推一把：只有进动频率与之匹配的质子被“推”起来，宏观磁化被扳倒、进动、再缓慢恢复，恢复过程中辐射出的微弱射频信号被线圈接收。不同组织（脂肪、水、肿瘤）里质子恢复的步调不同，信号对比就成了图像。为什么开场说“点名”而非“拍照”：磁共振的本质，是用共振这把音叉，从亿万个质子里把频率合拍的那一群请出来合唱。这里没有射线、没有电离，全部机关只是本章的两个事实——自旋只有两种取向，磁矩跟着自旋走。检查一个极端情形：把磁场关掉，$\Delta E\to0$，共振频率跌到零，射频脉冲什么也选不出来，整机哑火——图像对磁场的依赖是彻底的，这正是公式 $f\propto B$ 在 $B=0$ 端的表现。

顺带收掉一个可能的疑问：地磁场约 $5\times10^{-5}$~T，你的指南针感受不到量子劈裂，因为针是宏观物体，亿万个自旋与晶格搅在一起，平均磁化连续可变，且热扰动远大于单自旋的能级差——量子规则没有消失，只是被平均掩埋了。微观离散、宏观连续，又一次对上了。

\step{脑洞实验与挑战题}

\begin{thought}[磁星：把自旋物理推到宇宙的极端]
中子星是恒星坍缩后的残骸，半径只有十来千米，质量却比太阳还大。其中一类叫磁星，表面磁场高达 $10^{11}$~T——是磁共振仪的近千亿倍。用本章的尺子量一量那里的世界：质子在那样的磁场里，两种自旋取向的能量差对应约 $4$~GHz 的进动频率（微波波段）；电子自旋的能级劈裂还要再大三个数量级，落在软 X 射线波段。更直观地说，在磁星表面，原子里的电子被磁场约束成细针状，化学键被磁场改写，一块“普通物质”的结构都会变形。同一个 $\hbar/2$，在你家的冰箱贴上是几微牛·米的力矩，在磁星上是能扭曲原子本身的力量。物理规律没有换，换的只是参数走到了极端。
\end{thought}

\begin{puzzles}
\item 有人提出：“电子自旋的两种结果可以这样理解——每个电子内部早就写好了沿 $x$、$y$、$z$ 三个方向的答案（各为 $+$ 或 $-$），测量只是读出其一。”只用本章的串联斯特恩—盖拉赫实验结果，你能不能反驳这个模型？\textbf{提示}：如果“$z$ 方向的答案”一直写在电子身上，那么中间插入一次 $x$ 方向测量后，再测 $z$ 时结果应该仍然是第一次筛出的“上”；实验给了什么？再想想：这个反驳只排除了“$z$ 答案不被 $x$ 测量打扰”的版本，更顽强的“每个方向都有预定答案”的版本要等第~48~章的贝尔不等式来终审。
\item 一束银原子通过沿 $z$ 方向的装置，留下“上”通道；接着让它通过一台沿与 $z$ 夹角 $60^\circ$ 方向的装置。问两个通道的强度比是多少？如果之后再加一台转回 $z$ 方向的装置，仅用第一次留下的“上”原子做输入，你预期看到什么？\textbf{提示}：用 $P=\cos^2(\theta/2)$，$\theta=60^\circ$ 时 $\cos^2 30^\circ=\tfrac34$；最后一问先别急着套数字——经过中间那次测量后，原子的状态已经变了。
\item 经典陀螺仪靠角动量方向稳定来导航。现在设想用单个电子的自旋做“量子陀螺仪”，它能不能像经典陀螺那样给出一个连续可读的指向？为什么工程师实际上使用的是\textbf{大量}自旋的系综（如核磁共振陀螺仪）？\textbf{提示}：单个自旋单次测量只有两种结果；连续的指向信息藏在哪里？回忆“微观离散、平均连续”。
\item 把一块钕磁铁在酒精灯上烧热，再冷却，它的磁性明显变弱（请勿用强磁铁靠近火源太久，本题为思想题）。结合本章与热学知识解释：热运动对自旋大军做了什么？为什么冷却之后磁性不能完全自动恢复？\textbf{提示}：排队需要能量上划算，热运动会随机翻转自旋；冷却时自旋会重新排队，但相邻区域的队列方向未必一致——想想“畴”，即各自为政的小队列。
\end{puzzles}

下一章，我们把两个这样的自旋放在一起。你将会看到，量子力学里最著名的“鬼魅般的超距作用”——纠缠——恰恰是用本章的自旋当主角被贝尔不等式钉上实验审判台的；而量子计算机的所有魔力，也不过在布洛赫球上跳一支多人舞。
